% ------------------------------------------------------------
% ------------------------------------------------------------
  
% Insert the common KETCube AppNote Defines

% ------------------------------------------------------------
% ------------------------------------------------------------

\documentclass[twoside,a4paper,usenames,dvipsnames]{refart}
\usepackage[utf8x]{inputenc}
%\usepackage[czech]{babel}
\usepackage[pdftex]{graphicx}
\graphicspath{{resources/images/}}
\usepackage{caption}% for \captionof
\usepackage{mwe}% contains example-image

% Definition of shortcuts for authornames
%  Sort alphabetically by surname

\newcommand{\JB}{Jan Bělohoubek}
\newcommand{\JC}{Jiří Čengery}
\newcommand{\JF}{Jaroslav Freisleben}
\newcommand{\PK}{Petr Kašpar}
\newcommand{\MU}{Matin Úbl}
\newcommand{\KV}{Kryštof Vaněk}
\newcommand{\JZ}{Jan Záruba}


\usepackage[owncaptions]{vhistory}
\usepackage{hyperref}
%\usepackage[superscript,biblabel]{cite}
\usepackage{multirow}
\usepackage{wrapfig}
\usepackage{float}
\usepackage{siunitx}
\usepackage[table, x11names]{xcolor}
\usepackage{fancyvrb}
\usepackage{fontawesome}
\usepackage[many]{tcolorbox}
\usepackage{listings}
\usepackage{tikz}
\usetikzlibrary{shapes,positioning}
\usepackage{verbatimbox}

\usepackage{array, booktabs, boldline} %
\usepackage{mathtools}
\usepackage[normalem]{ulem}

\usepackage[T1]{fontenc}% Needed for \textquotedbl
\DeclareSIUnit[number-unit-product = {}]{\inchQ}{\textquotedbl}

\newsavebox{\tempbox}
\newlength{\tempheight}

\newcommand\ToDo[1]{\textcolor{red}{ToDo: #1}}

% --- Boxed Enviroments

\newcommand\docNote[1]{
\par
\addvspace{0.5cm}
\subsubsection*{~\hspace{0.2cm}\textcolor{gray}{\Huge \faCommentingO\,}}
\vspace{-1.5cm}
\begin{tcolorbox}[breakable,colback=white,colframe=gray,width=\dimexpr\textwidth+32mm\relax,enlarge left by=-32mm, title = Note]
\it \small #1
\end{tcolorbox}
}

\newcommand\docWarn[1]{
\par
\addvspace{0.5cm}
\subsubsection*{~\hspace{0.2cm}\textcolor{gray}{\Huge \faWarning\,}}
\vspace{-1.5cm}
\begin{tcolorbox}[breakable,colback=white,colframe=gray,width=\dimexpr\textwidth+32mm\relax,enlarge left by=-32mm, title = Warning]
\it #1
\end{tcolorbox}
}

\newcommand\docExample[1]{
\par
\addvspace{0.5cm}
\subsubsection*{~\hspace{0.2cm}\textcolor{gray}{\Huge \faCogs\,}}
\vspace{-1.5cm}
\begin{tcolorbox}[breakable,colback=white,colframe=gray,width=\dimexpr\textwidth+32mm\relax,enlarge left by=-32mm, title = Example]
#1
\end{tcolorbox}
}

\newenvironment{docCodeExample}
{
\par
\addvspace{0.5cm}
\subsubsection*{~\hspace{0.2cm}\textcolor{gray}{\Huge \faCogs\,}}
\vspace{-1.5cm}
\begin{tcolorbox}[breakable,colback=white,colframe=gray,width=\dimexpr\textwidth+32mm\relax,enlarge left by=-32mm, title = Example]
}
{
\end{tcolorbox}
}

\newenvironment{docCodeExampleTitled}[1]
{
\par
\addvspace{0.5cm}
\subsubsection*{~\hspace{0.2cm}\textcolor{gray}{\Huge \faCogs\,}}
\vspace{-1.5cm}
\begin{tcolorbox}[breakable,colback=white,colframe=gray,width=\dimexpr\textwidth+32mm\relax,enlarge left by=-32mm, title = Example: #1]
}
{
\end{tcolorbox}
}

% --- Filesystem Names 
\newcommand\docPath[1]{{\tt #1}}
\newcommand\docFileName[1]{{\tt #1}}

% --- Source Code Names 
\newcommand\docVarName[1]{{\tt #1}}
\newcommand\docFnName[1]{{\tt #1}}
\newcommand\docTypeName[1]{{\tt #1}}

% --- KETCube Terms
\newcommand\docKCModName[1]{{\it #1} module}
\newcommand\docKCCmdInline[1]{\colorbox{gray!30}{\tt #1}}
\newcommand\docKCCmd[1]{{\tt > #1}}

% --- NICE sections
\newcommand\niceSubSection[2]{
\subsection*{~\hspace{0.2cm}{\Huge #1}\\[-0.6cm]\phantom{x}~\hspace{1cm}~#2}
\vspace{-0.7cm}
}

%skryt barevny obdelnik kolem odkazu
\hypersetup{
    colorlinks=false,
    pdfborder={0 0 0},
}

% vhistory
\renewcommand{\vhhistoryname}{Revision History}
\renewcommand{\vhchangename}{Note}
\renewcommand{\vhversionname}{Revision}
\renewcommand{\vhdatename}{Date}
\renewcommand{\vhauthorname}{Author}
\renewcommand \vhAuthorColWidth{0.8\hsize}
\renewcommand \vhChangeColWidth{1.2\hsize}

\DeclareRobustCommand{\UWBLogo}{%
   \begin{wrapfigure}{l}{2.1cm}
    \vspace{-1.35cm}
    \includegraphics[width=2cm]{ZCU_logo.pdf}
   \end{wrapfigure}
}


% declare the path(s) where your graphic files are
% extract presenattion number to create include path for images
\ExplSyntaxOn
% Save a copy of \jobname
\tl_set:NV \NUMBER \c_sys_jobname_str
\regex_replace_once:nnN { [A-Za-z]*_[A-Za-z]*_ } { } \NUMBER
\ExplSyntaxOff

% declare the path(s) where your graphic files are
%\graphicspath{{resources/images/}{resources/appNotes/001/images/}}
\graphicspath{{resources/images/}{resources/appNotes/\NUMBER/images/}}
\graphicspath{{resources/images/}{resources/appNotes/003/images/}}

\pdfinfo
{
  /Title       (The KETCube Project AppNote)
  /Creator     (LaTeX)
  /Author      (The SmartCAMPUS Team)
}




\title{\UWBLogo KETCube AppNote 006:\\ KETCube Upstream Modules (\vhCurrentVersion)}

\author{Author: \vhListAllAuthorsLongWithAbbrev}
\date{Version \vhCurrentVersion\ from \vhCurrentDate}

% ------------------------------------------------------------
% ------------------------------------------------------------
  
% Insert the common KETCube AppNote Head

% ------------------------------------------------------------
% ------------------------------------------------------------

\begin{document}
\pagenumbering{roman} 

\titlepage
\maketitle

% ------------------------------------------------------------
% ------------------------------------------------------------
  
% BEGIN of the KETCube appNote Content

% ------------------------------------------------------------
% ------------------------------------------------------------


% ------------------------------------------------------------
% ------------------------------------------------------------
  
% BEGIN of the KETCube appNote Content

% ------------------------------------------------------------
% ------------------------------------------------------------

\section*{About this Document}
{\it KETCube} \cite{ZCU:KETCube:05-2018} is the prototyping and demo platform developed at the Department of Materials and Technology (KET), University of West Bohemia in Pilsen. 


This document describes KETCube upstream modules. 

\setcounter{tocdepth}{2}
\tableofcontents
\clearpage

\listoffigures
\listoftables
\begin{versionhistory}
  \vhEntry{0.2.0*}{16.09.2019}{JB}{Initial version}
\end{versionhistory}
% history table ... do not number
\setcounter{table}{0}

\clearpage 
\pagenumbering{arabic} 
\pagestyle{headings} 

\clearpage

\section{Module Activation}\label{sec:activation}
KETCube is a modular platform, where number of modules is included in the platform release. Any module can be individually enabled or disabled based on user needs and extension boards connected to the KETCube {\it Main Board}. For details, please refer to the KETCube Datasheet \cite{ZCU:KETCube:05-2018}.

To display list of available modules, open the KETCube Terminal and use \docKCCmdInline{list} command. Commands \docKCCmdInline{enable}/\docKCCmdInline{disable} are used to turn ON/OFF modules (e.g. \docKCCmdInline{enable HDC1080}). When module is enabled, it starts to perform defined operation (e.g. measure RH and Temperature).

Debug messages are useful when device is initially configured or in case, that unexpected behavior occurred.

Debug messages are printed to serial terminal interface. They can be configured on the per-module basis by setting the {\it severity level} to the selected module or to the KETCube core.

Four severity levels are defined:
\begin{itemize}
  \item[0 --] NONE -- no messages enabled
  \item[1 --] ERROR -- only error messages enabled (default severity)
  \item[2 --] INFO -- error and informational messages enabled
  \item[3 --] DEBUG -- previous message groups and debug information are enabled
\end{itemize}

The second (optional) parameter of the “enable” command is used to configure the {\it severity level} of the selected module (e.g. \docKCCmdInline{enable HDC1080 3} enables debug messages for HDC1080 module). The command \docKCCmdInline{set core severity} is used to configure \docKCCmdInline{severity level} of the KETCube core and command \docKCCmdInline{set driver severity} is used to configure {\it severity level} of the KETCube low-level drivers.

\docNote{Debug messages may cause a situation, where it may be difficult to write commands (terminal echo is foggy due to lot of debug messages produced by KETCube)}.

The commands \docKCCmdInline{show(r)}/\docKCCmdInline{set(r)} are used to show/set KETCube (module) settings. The commands are organized in a tree-like manner: the \docKCCmdInline{show(r)}/\docKCCmdInline{set(r)} command is the root command, it is followed by the module name and by the module variable, which is to be set. 

\begin{docCodeExample}
\begin{verbatim}
>> show LoRa devEUI
[root command] [module name] [variable name]
\end{verbatim}
\end{docCodeExample}

\docNote{The persistent parameters are saved into on-chip EEPROM, the running configuration is stored in RAM.}

% ====================
\clearpage
\section{KETCube Actuation Modules}
\docNote{No actuation modules are currently part of the KETCube upstream.}
  
% ====================
\clearpage
\section{KETCube Communication Modules}

\subsection{LoRa Module}
\docKCModName{LoRa} allows to use KETCube as a LoRaWAN node device. 

The \docKCModName{LoRa} can be enabled the same way as any other module -- see Section \ref{sec:activation}. Most of the LoRaWAN parameters can be set by using the KETCube Terminal interface.

Write command \docKCCmdInline{show LoRa} or \docKCCmdInline{set LoRa} to display commands related to the \docKCModName{LoRa}:
\begin{docCodeExample}
\begin{verbatim}
>> show LoRa
appEUI              LoRa application EUI
appKey              LoRa application key
appSKey             LoRa app session Key
devAddr             LoRa device address
devEUIBoard         Board (boardID-based) LoRa device EUI
devEUICustom        Custom LoRa device EUI
enableABP           Enable ABP
enableCustomDevEUI  Custom (user-defined) LoRa device EUI
nwkSKey             LoRa network session Key

>> set LoRa
appEUI              LoRa application EUI
appKey              LoRa application key
appSKey             LoRa app session Key
devAddr             LoRa device address
devEUICustom        Custom LoRa device EUI
enableABP           Enable ABP
enableCustomDevEUI  Custom (user-defined) LoRa device EUI
nwkSKey             LoRa network session Key
\end{verbatim}
\end{docCodeExample}

\subsubsection{LoRaWAN Node Activation}
To select or check OTAA/ABP mode settings, the following commands are used: \docKCCmdInline{set LoRa enableABP {0|1}}/\docKCCmdInline{show LoRa enableABP}. The “set” command parameter is boolean {\tt 0} or {\tt 1}. OTAA is the default mode.

\subsection{LoRaWAN devEUI}
User-defined devEUI (\docKCCmdInline{set LoRa enableCustomDevEUI 1} or manufacturer's devEUI (\docKCCmdInline{set LoRa enableCustomDevEUI 0}) can be used. The selected option can be checked via \docKCCmdInline{show LoRa enableCustomDevEUI} command. The user-defined devEUI can be configured by \docKCCmdInline{set LoRa devEUICustom} command. Manufacturer's devEUI is used by default.

\clearpage

\begin{docCodeExampleTitled}{Set user-defined devEUI}
\begin{verbatim}
>> set LoRa devEUICustom 1122334455667788
>> set LoRa enableCustomDevEUI 1
>> reload

>> show LoRa devEUICustom

Custom devEUI is displayed ...
\end{verbatim}
\end{docCodeExampleTitled}

\begin{docCodeExampleTitled}{Use Manufacturer's devEUI}
\begin{verbatim}
>> set LoRa enableCustomDevEUI 0
>> reload

>> show LoRa devEUIBoard

Board devEUI is displayed ...
\end{verbatim}
\end{docCodeExampleTitled}

\subsubsection{LoRaWAN OTAA Parameters}
The LoRaWAN appEUI and appKey can be set by using \docKCCmdInline{LoRa appEUI} and \docKCCmdInline{set LoRa appKey} commands and checked by using \docKCCmdInline{show LoRa appEUI} and \docKCCmdInline{show LoRa appKey} commands.
  
\begin{docCodeExampleTitled}{Set OTAA parameters}
\begin{verbatim}
>> set LoRa appEUI 1122334455667788
>> set LoRa appKey 11223344556677881122334455667788
\end{verbatim}
\end{docCodeExampleTitled}
  
\subsubsection{LoRaWAN ABP Parameters}
The LoRaWAN appSKey and nwkSKey can be set by using \docKCCmdInline{set LoRa appSKey} and \docKCCmdInline{set LoRa nwkSKey} commands and checked by using \docKCCmdInline{show LoRa appSKey} and \docKCCmdInline{show LoRa nwkSKey} commands. The static device address i configured via  \docKCCmdInline{set LoRa devAddr} command and checked via  \docKCCmdInline{show LoRa devAddr} command.
  
\begin{docCodeExampleTitled}{Set ABP parameters}
\begin{verbatim}
>> set LoRa devAddr DEADBEEF
>> set LoRa appSKey 1122334455667788
>> set LoRa nwkSKey 11223344556677881122334455667788
\end{verbatim}
\end{docCodeExampleTitled}

\subsubsection{LoRaWAN Period}
The period of data transmission is given by the {\it basePeriod} -- see the KETCube Datasheet \cite{ZCU:KETCube:05-2018}.

\subsubsection{LoRaWAN Ports}

  The periodical messages are transmitted by using LoRaWAN port 2. Asynchronous messages (particular modules can indicate errors asynchronously, push-button event) are transmitted by using LoRaWAN port 3 if and only if the \docKCModName{AsyncTx} is enabled.

\clearpage
\subsection{StarNetNode Module}

\clearpage
\subsection{StarNetConcentrator Module}

% ====================
\clearpage
\section{KETCube Sensing Modules}

\subsection{HDC1080 Module}

The HDC1080 is on-board {\it Relative Humidity} and {\it Temperature} sensor allowing RH+T measurement. It can be enabled the same way as any other module -- see Section \ref{sec:activation}.

\subsubsection{HDC1080 Module Configuration}
  The \docKCModName{HDC1080} only allows to configure severity level, other configuration options are not available.
  
\begin{docCodeExample}
\begin{verbatim}
>> enable HDC1080 2

   ...
>> reload
Executing command: reload

Performing system reset and reloading KETCube configuration ...  
\end{verbatim}
\end{docCodeExample}
  
  Once reloaded, the \docKCModName{HDC1080} measures RH+T periodically. 
  
  For debugging purposes, it may be useful to watch RH+T output in console. The output for the severity level 2 (INFO) is shown below:

\begin{docCodeExample}
\begin{verbatim}
>> HDC1080 :: Temperature: 25, RH: 57
\end{verbatim}
\end{docCodeExample}

\subsubsection{HDC1080 Module Output}
  The \docKCModName{HDC1080} output consist of 4 bytes:
  
\begin{tikzpicture}
 \node[draw,fill=SeaGreen3,minimum height=1cm,minimum width = 2cm,xshift=0cm](10/5){$T_{MSB}$};
 \node[draw,fill=SeaGreen3,minimum height=1cm,minimum width = 2cm,xshift=2cm](10/5){$T_{LSB}$};
 \node[draw,fill=white,minimum height=1cm,minimum width = 2cm,xshift=4cm](10/5){$RH_{MSB}$};
 \node[draw,fill=white,minimum height=1cm,minimum width = 2cm,xshift=6cm](10/5){$RH_{LSB}$};
\end{tikzpicture}

  Temperature is encoded in BigEndian as unsigned integer in additive code. The value units are degrees of Celsius.
  
  To compute actual temperature, use the following equation:
  
  $T = ((T_{MSB} << 8 ~|~ T_{LSB}) - 10000) / 10 ~[°C]$
  
  Relative Humidity is encoded in BigEndian format as unsigned integer in additive code. The value units are \% RH.
  
  To compute actual RH, use the following equation:
  
  $RH = (RH_{MSB} << 8 ~|~ RH_{LSB}) / 10 ~[\%]$
  
  An error for both temperature and humidity is indicated by out-of-the range values from the unsigned word (16-bit) range -- see ranges in Section \ref{spec:OC}.

% ------------------
\clearpage
\subsection{BatMeas Module}
  
The \docKCModName{BatMeas} module allows to monitor KETCube battery (if connected). 

The \docKCModName{BatMeas} can be enabled the same way as any other module -- see Section \ref{sec:activation}.

\subsubsection{BatMeas Module Configuration}
  Use the commands {\tt list} and {\tt bat} to select connected battery.
  
  The sample configuration procedure:
  
\begin{docCodeExample}
\begin{verbatim}
>> show batMeas list
Executing command: list
Available batteries:
0)       CR2032 (Up to 560mAh battery)
1)       LS33600 (15 Ah battery)
>> set batMeas bat 0
Executing command: bat

Write success!

>> enable BatMeas
Executing command: enable

Setting module BatMeas:     success!
>> reload
Executing command: reload

Performing system reset and reloading KETCube configuration ...
\end{verbatim}
\end{docCodeExample}

  Once reloaded, the \docKCModName{BatMeas} measures battery voltage periodically (voltage at KETCube 3V3 PIN). 
  
For debugging purposes, it may be useful to watch the \docKCModName{BatMeas} output in console. The output for the severity level 2 (INFO) is shown below:

\begin{docCodeExample}
\begin{verbatim}
>> batLevel :: 0
\end{verbatim}
\end{docCodeExample}

\subsubsection{BatMeas Module Output}
  The \docKCModName{BatMeas} output consist of a single byte only.
  
  The \docKCModName{BatMeas} output value (B) meaning is as follows:
  
  \begin{itemize}
    \item 0x00: battery is (almost) discharged
    \item 0x01 -- 0xFE: battery voltage scaled to 1 -- 254
    \item 0xFF: RFU
  \end{itemize}
  
  To compute actual battery voltage, use the following equation:
  
  $V = (B/254 \cdot 0.7) + V_{min} ~[mV]$
  
  $V_{min}$ is the lowest voltage value -- battery is considered as discharged. $V_{min}$ and $V_{nom}$ (nominal voltage) values for supported batteries are in the following table:
  
  \begin{table*}[!ht]
    %\hspace*{-4cm}
    \begin{tabular}{| p{3.5cm} | p{2cm} | p{2cm} |}
        \hline
        \rowcolor{SeaGreen3!30!} {\bf Battery} & $V_{min}$ [V] & $V_{nom}$ [V] \\
        \hline
        \hline
        CR2032 & 2.9 & 3.3 \\
        \hline
        LS33600 & 2.9 & 3.6 \\
        \hline
    \end{tabular}
    \addcontentsline{lot}{table}{Supported Batteries}
    \label{tab:spec:AMR}
   \end{table*}
  
% ------------------------------------------------------------
% ------------------------------------------------------------
  
% END of the KETCube appNote Content

% ------------------------------------------------------------
% ------------------------------------------------------------



% ------------------------------------------------------------
% ------------------------------------------------------------
  
% Insert the common KETCube appNote Tail

% ------------------------------------------------------------
% ------------------------------------------------------------

% ------------------------------------------------------------
% ------------------------------------------------------------
  
% END of the KETCube appNote Content

% ------------------------------------------------------------
% ------------------------------------------------------------


\clearpage
\bibliographystyle{IEEEtran}
\bibliography{IEEEabrv,resources/sources}

%
% Include this license into all KETCube-related documentation
%
%

\clearpage

~

\vfill

\section*{Important Notice}

\subsection*{Copyright}
\copyright ~2018 University of West Bohemia in Pilsen\\
All rights reserved.

\subsection*{Developed by}
The SmartCAMPUS Team\\
Department of Technologies and Measurement\\
Faculty of Electrical Engineering\\
www.smartcampus.cz/en $\mid$ www.zcu.cz/en

\subsection*{License\footnotemark}
\footnotetext{{\it University of Illinois/NCSA Open Source License} (\url{https://opensource.org/licenses/NCSA}) -- similar to {\it Modified BSD License}}

Permission is hereby granted, free of charge, to any person obtaining a copy of this software and associated documentation files (the “Software”), to deal with the Software without restriction, including without limitation the rights to use, copy, modify, merge, publish, distribute, sublicense, and/or sell copies of the Software, and to permit persons to whom the Software is furnished to do so, subject to the following conditions:

\begin{itemize}
    \item[--] Redistributions of source code must retain the above copyright notice, this list of conditions and the following disclaimers.
    \item[--] Redistributions in binary form must reproduce the above copyright notice, this list of conditions and the following disclaimers in the documentation and/or other materials provided with the distribution.
    \item[--] Neither the names of The SmartCAMPUS Team, Department of Technologies and Measurement and Faculty of Electrical Engineering University of West Bohemia in Pilsen, nor the names of its contributors may be used to endorse or promote products derived from this Software without specific prior written permission. 
\end{itemize}
THE SOFTWARE IS PROVIDED “AS IS”, WITHOUT WARRANTY OF ANY KIND, EXPRESS OR IMPLIED, INCLUDING BUT NOT LIMITED TO THE WARRANTIES OF MERCHANTABILITY, FITNESS FOR A PARTICULAR PURPOSE AND NONINFRINGEMENT. IN NO EVENT SHALL THE CONTRIBUTORS OR COPYRIGHT HOLDERS BE LIABLE FOR ANY CLAIM, DAMAGES OR OTHER LIABILITY, WHETHER IN AN ACTION OF CONTRACT, TORT OR OTHERWISE, ARISING FROM, OUT OF OR IN CONNECTION WITH THE SOFTWARE OR THE USE OR OTHER DEALINGS WITH THE SOFTWARE. 

\end{document}


\end{document}


